% ---------------------------------------------------------------------
% RECUPERO (file R1, R2) - diagramma di flusso con ciclo (senza suggerimenti) + tabella
%   fila R1: ascensore, persone dello stesso peso -> persone salite e peso totale
%            (condizione T+P <= M, come lo scaffale della fila C)
%   fila R2: robot a batteria, compiti con consumo fisso -> compiti svolti e carica rimasta
%            (condizione B >= E, come il magazzino della fila B)
%   Prova: R1 M=300, P=70 -> 4 iterazioni;  R2 B=100, E=22 -> 4 iterazioni
% ---------------------------------------------------------------------
\begin{esercizio}[14]
\RR{%
Su un ascensore salgono persone, una alla volta, tutte dello stesso peso. Disegna il diagramma di
flusso di un algoritmo che conta quante persone possono salire.\\
\textbf{Dati da leggere:} $M$ = peso massimo che l'ascensore può portare (in kg); $P$ = peso di una
persona (in kg).\\
\textbf{Altre variabili:} $T$ = peso già sull'ascensore; $N$ = numero di persone salite.\\
Le persone salgono finché ne può salire ancora una senza superare il peso massimo $M$.
Alla fine visualizza $N$ e $T$.\\
\textbf{Prova:} con $M = 300$ e $P = 70$ completa la tabella con i valori dopo ogni iterazione del ciclo.
\begin{center}\renewcommand{\arraystretch}{1.25}
\begin{tabular}{|c|p{3cm}|p{3cm}|}\hline
  Iterazione & $T$ & $N$\\\hline
  Inizio & &\\\hline
  1 & &\\\hline
  2 & &\\\hline
  3 & &\\\hline
  4 & &\\\hline
  5 & &\\\hline
  6 & &\\\hline
\end{tabular}\end{center}}{%
Un robot svolge compiti uno dopo l'altro e ogni compito consuma sempre la stessa parte della
batteria. Disegna il diagramma di flusso di un algoritmo che conta quanti compiti riesce a svolgere.\\
\textbf{Dati da leggere:} $B$ = carica della batteria (in \%); $E$ = carica consumata da un compito (in \%).\\
\textbf{Altre variabili:} $K$ = numero di compiti svolti.\\
Il robot svolge compiti finché la batteria ha abbastanza carica per un compito intero.
Alla fine visualizza $K$ e la carica rimasta $B$.\\
\textbf{Prova:} con $B = 100$ e $E = 22$ completa la tabella con i valori dopo ogni iterazione del ciclo.
\begin{center}\renewcommand{\arraystretch}{1.25}
\begin{tabular}{|c|p{3cm}|p{3cm}|}\hline
  Iterazione & $B$ & $K$\\\hline
  Inizio & &\\\hline
  1 & &\\\hline
  2 & &\\\hline
  3 & &\\\hline
  4 & &\\\hline
  5 & &\\\hline
  6 & &\\\hline
\end{tabular}\end{center}}

\begin{solution}
\begin{center}
\RR{%
\begin{tikzpicture}[x=1cm,y=1cm]
  \node[fcterm] (ini)  at (0,0) {Inizio};
  \node[fcio]   (in)   at (0,-1.1) {Leggi $M$, $P$};
  \node[fcinit] (ini0) at (0,-2.2) {$T \gets 0$, \ $N \gets 0$};
  \node[fcdec]  (d1)   at (0,-3.75) {$T + P \leq M$};
  \node[fcproc] (agg)  at (0,-5.3) {$T \gets T + P$};
  \node[fcproc] (inc)  at (0,-6.4) {$N \gets N + 1$};
  \node[fcio]   (out)  at (5.0,-3.75) {Scrivi $N$, $T$};
  \node[fcterm] (fin)  at (5.0,-4.95) {Fine};
  \draw[fcfreccia] (ini) -- (in);
  \draw[fcfreccia] (in) -- (ini0);
  \draw[fcfreccia] (ini0) -- (d1);
  \draw[fcfreccia] (d1) -- node[fcsi,right] {Sì} (agg);
  \draw[fcfreccia] (agg) -- (inc);
  \draw[fcfreccia] (inc.south) -- ++(0,-0.5) -| (-2.4,-2.95) -- (0,-2.95);
  \draw[fcfreccia] (d1) -- node[fcno,above] {No} (out);
  \draw[fcfreccia] (out) -- (fin);
\end{tikzpicture}}{%
\begin{tikzpicture}[x=1cm,y=1cm]
  \node[fcterm] (ini)  at (0,0) {Inizio};
  \node[fcio]   (in)   at (0,-1.1) {Leggi $B$, $E$};
  \node[fcinit] (ini0) at (0,-2.2) {$K \gets 0$};
  \node[fcdec]  (d1)   at (0,-3.75) {$B \geq E$};
  \node[fcproc] (agg)  at (0,-5.3) {$B \gets B - E$};
  \node[fcproc] (inc)  at (0,-6.4) {$K \gets K + 1$};
  \node[fcio]   (out)  at (5.0,-3.75) {Scrivi $K$, $B$};
  \node[fcterm] (fin)  at (5.0,-4.95) {Fine};
  \draw[fcfreccia] (ini) -- (in);
  \draw[fcfreccia] (in) -- (ini0);
  \draw[fcfreccia] (ini0) -- (d1);
  \draw[fcfreccia] (d1) -- node[fcsi,right] {Sì} (agg);
  \draw[fcfreccia] (agg) -- (inc);
  \draw[fcfreccia] (inc.south) -- ++(0,-0.5) -| (-2.4,-2.95) -- (0,-2.95);
  \draw[fcfreccia] (d1) -- node[fcno,above] {No} (out);
  \draw[fcfreccia] (out) -- (fin);
\end{tikzpicture}}
\end{center}
\RR{%
  Ogni iterazione del ciclo è una persona che sale: $N$ conta le iterazioni. Prima di far salire una
  persona si controlla che il peso con la persona in più ($T+P$) non superi $M$; alla fine $T \leq M$.\\[4pt]
  \textbf{Prova} ($M = 300$, $P = 70$):
\begin{center}\renewcommand{\arraystretch}{1.1}
\begin{tabular}{|c|>{\centering\arraybackslash}p{2.5cm}|>{\centering\arraybackslash}p{2.5cm}|}\hline
  Iterazione & $T$ & $N$\\\hline
  Inizio & 0 & 0\\\hline
  1 & 70 & 1\\\hline
  2 & 140 & 2\\\hline
  3 & 210 & 3\\\hline
  4 & 280 & 4\\\hline
\end{tabular}\end{center}  Dopo la 4ª iterazione $T = 280$ e $280 + 70 = 350 \leq 300$ è falso: il ciclo finisce. Si visualizza $N = 4$ e $T = 280$.}{%
  Ogni iterazione del ciclo è un compito svolto: $K$ conta le iterazioni. Quando il ciclo finisce,
  nella batteria resta meno carica di quella che serve per un compito.\\[4pt]
  \textbf{Prova} ($B = 100$, $E = 22$):
\begin{center}\renewcommand{\arraystretch}{1.1}
\begin{tabular}{|c|>{\centering\arraybackslash}p{2.5cm}|>{\centering\arraybackslash}p{2.5cm}|}\hline
  Iterazione & $B$ & $K$\\\hline
  Inizio & 100 & 0\\\hline
  1 & 78 & 1\\\hline
  2 & 56 & 2\\\hline
  3 & 34 & 3\\\hline
  4 & 12 & 4\\\hline
\end{tabular}\end{center}  Dopo la 4ª iterazione $B = 12$ e $12 \geq 22$ è falso: il ciclo finisce. Si visualizza $K = 4$ e $B = 12$.}
  \criterioRR
  {Inizio/Fine e simboli corretti (inizializzazione nell'esagono): 1 pt; lettura di $M$ e $P$: 1 pt;
    inizializzazione $T\gets 0$, $N\gets 0$: 2 pt; condizione del ciclo $T+P\leq M$ ricavata da ``ne può
    salire ancora una'', con rami Sì/No giusti: 3 pt ($T<M$ o $T\leq M$, che fanno superare $M$: 1 pt);
    aggiornamento $T\gets T+P$: 2 pt; contatore $N\gets N+1$: 1 pt; ritorno alla condizione: 1 pt; uscita
    di $N$ e $T$ dopo il ciclo: 1 pt. Ciclo con condizione in fondo corretto se gestisce il caso $P>M$:
    stessi punti. Tabella della prova: 2 pt ($-0{,}5$ per ogni riga sbagliata; se il diagramma è sbagliato, vale la tabella coerente con il diagramma disegnato). Valori attesi: $T$ = 0, 70, 140, 210, 280; $N$ = 0, 1, 2, 3, 4.}
  {Inizio/Fine e simboli corretti (inizializzazione nell'esagono): 1 pt; lettura di $B$ e $E$: 1 pt;
    inizializzazione $K\gets 0$ (e nessuna inizializzazione di $B$ dopo la lettura): 2 pt; condizione del
    ciclo $B\geq E$ ricavata da ``abbastanza carica per un compito intero'', con rami Sì/No giusti:
    3 pt ($B>E$, che perde il compito quando $B=E$: 1 pt); aggiornamento $B\gets B-E$: 2 pt; contatore
    $K\gets K+1$: 1 pt; ritorno alla condizione: 1 pt; uscita di $K$ e $B$ dopo il ciclo: 1 pt. Ciclo con
    condizione in fondo corretto se gestisce il caso $B<E$ iniziale: stessi punti. Tabella della prova: 2 pt ($-0{,}5$ per ogni riga sbagliata; se il diagramma è sbagliato, vale la tabella coerente con il diagramma disegnato). Valori attesi: $B$ = 100, 78, 56, 34, 12; $K$ = 0, 1, 2, 3, 4.}
\end{solution}
\end{esercizio}
